% Part: first-order-logic
% Chapter: introduction
% Section: first-order-logic

\documentclass[../../../include/open-logic-section]{subfiles}

\begin{document}

\olfileid{fol}{int}{fol}

\olsection{Eerste-ordelogica}

U kent de eerste-ordelogica waarschijnlijk uit een eerste inleiding tot
de formele logica.\footnote{We nemen dit min of meer aan!{} Zo niet, dan
kunt u een elementairder leerboek raadplegen, zoals \emph{forall x}
\citep{Magnus2021}.} Mogelijk kent u haar als ``kwantorlogica'' of
``predicatenlogica''. De eerste-ordelogica is om te beginnen een formele
taal. Zij heeft dus een bepaalde woordenschat en haar uitdrukkingen zijn
rijen symbolen uit die woordenschat. Niet iedere rij is echter
toegestaan. Er zijn verschillende soorten toegestane uitdrukkingen:
termen, !!{formula}s en !!{sentence}s. Onze belangstelling gaat vooral
uit naar !!{sentence}s van de eerste-ordelogica. Zij vormen een formele
tegenhanger van zinnen in het Nederlands, en we kunnen er de vragen over
stellen waarin een logicus doorgaans geïnteresseerd is. Bijvoorbeeld:
\begin{itemize}
    \item Volgt $!B$ logisch uit~$!A$?
    \item Is $!A$ logisch waar, logisch onwaar of
contingent?
    \item Zijn $!A$ en $!B$ equivalent?
\end{itemize}

Deze vragen betreffen primair de ``betekenis'' van !!{sentence}s van de
eerste-ordelogica. Een filosoof analyseert de vraag of $!B$ logisch uit
$!A$ volgt bijvoorbeeld als volgt: bestaat er een geval waarin $!A$
waar en $!B$ onwaar is ($!B$ volgt niet uit~$!A$), of maakt elk geval
dat $!A$ waar maakt ook $!B$ waar ($!B$ volgt wel uit~$!A$)? We hebben
echter nog niet bepaald wat een ``geval'' is; dat is de taak van de
\emph{semantiek}. De semantiek van de eerste-ordelogica geeft een
wiskundig precies model van het intuïtieve filosofische begrip ``geval''
en, even belangrijk, van wat het betekent dat !!a{sentence}~$!A$
\emph{waar is in} een geval. Het wiskundig precieze model dat we zullen
ontwikkelen heet !!a{structure}. De relatie die ``waar in'' preciseert,
heet de relatie van \emph{vervulling}. We zullen dus definiëren wat
``$!A$ wordt vervuld in~$\Struct{M}$'' betekent (in symbolen:
$\Sat{M}{!A}$), voor !!{sentence}s~$!A$ en !!{structure}s~$\Struct{M}$.
Daarna kunnen we ook andere semantische termen, zoals ``volgt uit'' en
``is logisch waar'', precies definiëren. Met die definities kunnen we
wederom wiskundig precies
vaststellen of bijvoorbeeld $\lforall[x][(!A(x) \lif !B(x)),
\lexists[x][!A(x)] \Entails \lexists[x][!B(x)]]$. Het antwoord is
uiteraard ``ja''. Wie al heeft geleerd Nederlandse zinnen in
eerste-ordelogica te symboliseren, herkent hierin bijvoorbeeld de
symbolisering van: ``Alle mieren zijn insecten; er zijn mieren; dus zijn
er insecten.'' Dit argument is duidelijk geldig. Ons wiskundige model
van ``volgt uit'' voor de formele taal moet daarom hetzelfde antwoord
geven.

Een ander onderwerp dat u waarschijnlijk uit een eerste inleiding tot
de formele logica kent, zijn \emph{!!{derivation}s}. In een eerste cursus
formele logica hebt u vermoedelijk geoefend zulke !!{derivation}s te
vinden, wellicht zelfs !!a{derivation} die de bovenstaande
gevolgtrekking aantoont. Er bestaan veel verschillende vormen van
!!{derivation}s. Mogelijk hebt u gewerkt met ``natuurlijke deductie'' of
``semantische tableaus'', maar er zijn vele andere. Systemen voor formele
!!{derivation} verschaffen hulpmiddelen om de bovenstaande vragen van logici
te beantwoorden. Zo \emph{verifieert} een !!{derivation} in natuurlijke
deductie waarin $\lforall[x][(!A(x) \lif !B(x))$ en
$\lexists[x][!A(x)]$ premissen zijn en $\lexists[x][!B(x)]]$ de conclusie
(laatste regel) is, dat $\lexists[x][!B(x)]$ logisch volgt uit
$\lforall[x][(!A(x) \lif !B(x))]$ en $\lexists[x][!A(x)]$.

Waarom is dat zo? Op het eerste gezicht staan systemen voor formele !!{derivation}
los van de semantiek: een formele !!{derivation} rangschikt slechts
symbolen volgens bepaalde regels en vermeldt geen ``gevallen'' of ``waar
in''. Het verband tussen systemen voor formele !!{derivation} en semantiek moet door
metalogisch onderzoek worden vastgesteld. Daarvoor is bijvoorbeeld een
wiskundig bewijs nodig dat een formele !!{derivation} van
$\lexists[x][!B(x)]$ uit de premissen
$\lforall[x][(!A(x) \lif !B(x))]$ en $\lexists[x][!A(x)]$ mogelijk is dan
en slechts dan als $\lforall[x][(!A(x) \lif !B(x))$ en
$\lexists[x][!A(x)]$ samen $\lexists[x][!B(x)]]$ impliceren. Voordat we
dit kunnen bewijzen, is echter veel zorgvuldig werk nodig om de
definities van syntaxis en semantiek correct te formuleren.

\end{document}
